\documentclass[lettersize,journal]{IEEEtran}

% ── Core math ────────────────────────────────────────────────────────────────
\usepackage{amsmath,amsfonts,amssymb}
\usepackage{amsthm}
\usepackage{bm}
\usepackage{mathtools}

% ── Theorem environments ─────────────────────────────────────────────────────
\newtheorem{theorem}{Theorem}
\newtheorem{lemma}{Lemma}
\newtheorem{proposition}{Proposition}
\usepackage{subcaption}
% ── Algorithms ───────────────────────────────────────────────────────────────
\usepackage{algorithm}
\usepackage{algpseudocode}

% ── Tables ───────────────────────────────────────────────────────────────────
\usepackage{array}
\usepackage{booktabs}
\usepackage{tabularx}
\usepackage{multirow}

% ── Figures ──────────────────────────────────────────────────────────────────

\usepackage{graphicx}
\usepackage{stfloats}
\usepackage[utf8]{inputenc}
\usepackage{textcomp}
\usepackage{url}
\usepackage{verbatim}
\usepackage{pifont}
\newcommand{\cmark}{\ding{51}}
\newcommand{\xmark}{\ding{55}}
\usepackage{bbm}
\usepackage{enumitem}
\usepackage{balance}
\usepackage{comment}
\usepackage[dvipsnames]{xcolor}
\usepackage[switch]{lineno}
\usepackage{cite}
\usepackage[colorlinks=true,
            linkcolor=blue,
            citecolor=blue,
            urlcolor=blue]{hyperref}

\DeclareMathOperator*{\argmin}{argmin}
\DeclarePairedDelimiter{\nint}\lfloor\rceil
\DeclarePairedDelimiter{\abs}\lvert\rvert

% ── Helper: best-in-column highlighting and metric arrows ────────────────────
\newcommand{\best}[1]{\textbf{#1}}
\newcommand{\up}{$\uparrow$}
\newcommand{\down}{$\downarrow$}
\newcommand{\nd}{\textemdash}            % not defined / not applicable

% ── Helper: a clearly-marked figure placeholder box (compiles w/o assets) ────
% Usage: \figph{height}{caption text shown inside the box}
\newcommand{\figph}[2]{%
  \setlength{\fboxsep}{4pt}%
  \fcolorbox{black}{gray!10}{%
    \begin{minipage}[c][#1][c]{0.96\linewidth}%
      \centering\itshape\small [\,Figure placeholder\,]\\[2pt] #2%
    \end{minipage}}%
}

\hyphenation{op-tical net-works semi-conduc-tor IEEE-Xplore trans-form-er}

% =============================================================================
\begin{document}
% =============================================================================

\title{A Unified Framework for Exact Relevance Propagation in Vision Transformers}

\author{Sathiyamohan Nishankar,~\IEEEmembership{Student Member,~IEEE}
  \IEEEcompsocitemizethanks{%
    \IEEEcompsocthanksitem
    Nishankar Sathiyamohan is with the Department of Computer Engineering,
    University of Peradeniya, Sri Lanka.
    E-mail: e17230@eng.pdn.ac.lk}%
  \thanks{Manuscript received \today.}}

\markboth{IEEE Transactions on Pattern Analysis and Machine Intelligence}%
         {Sathiyamohan: A Unified Framework for Exact Relevance Propagation in Vision Transformers}

\maketitle

% =============================================================================
\begin{abstract}
As Vision Transformers (ViTs) evolve from isotropic patch models to hierarchical and hybrid architectures (e.g., Swin, PVT, EfficientViT), existing Layer-wise Relevance Propagation (LRP) methods fundamentally break down. Spatial reduction and patch merging operations violate standard LRP conservation rules, forcing practitioners to fall back on heuristic approaches like Grad-CAM, which fail to capture fine-grained object structure and collapse on linear cross-covariance attention. In this paper, we introduce Hierarchical LRP (HiLRP), a framework that extends conservation-based relevance propagation to flat, hierarchical, and hybrid Vision Transformers. We derive flow-conserving LRP rules for spatial reduction and patch merging as corollaries of a single proposition, prove them conservative and conditionally equivariant, and verify both at machine precision. HiLRP is the only method that yields conservation-valid relevance on hierarchical and hybrid transformers, where the naive extension of prior attention-aware LRP fails; it attains the strongest and most consistent localization (0.980 on the Pointing Game for PVT, 0.960 for EfficientViT) and the highest agreement with axiomatic Shapley values. We further expose the failure of proxy faithfulness metrics on robust ViTs, establish Segment-Shapley Agreement as a discriminating evaluation standard, and extend the same machinery to explain label-free self-supervised objectives. HiLRP enables novel XAI capabilities, including stage-localized relevance tracking, providing unprecedented insight into the hierarchical feature construction of modern Vision Transformers.
\end{abstract}

\begin{IEEEkeywords}
Explainable AI, Vision Transformers, Hierarchical LRP, Relevance Propagation,
Faithfulness, Localization, Shapley Agreement.
\end{IEEEkeywords}

% =============================================================================
\section{Introduction}
\label{sec:introduction}
\IEEEPARstart{V}{ision} Transformers (ViTs)~\cite{dosovitskiy2021vit} and their hierarchical successors~\cite{liu2021swin,wang2021pvt,cai2023efficientvit} have displaced convolutional neural networks (CNNs) as the default backbones for image understanding. As these models enter high-stakes applications, attribution methods that assign an importance score to every input pixel for a given prediction remain the primary tool for post-hoc interpretation. 

Historically, Layer-wise Relevance Propagation (LRP)~\cite{bach2015lrp} and its transformer variants, the attention-attribution of Chefer et~al.~\cite{chefer2021transformer} and the attention-aware AttnLRP~\cite{achtibat2024attnlrp}, have provided mathematically grounded explanations by ensuring that the sum of relevance scores approximates the model's output logit. However, a critical gap has emerged: existing LRP rules are defined only for isotropic transformers and are incompatible with modern hierarchical ViT architectures. While isotropic ViTs maintain a constant token sequence length, state-of-the-art models employ hierarchical designs featuring patch merging, spatial reduction, and linear cross-covariance attention. When standard LRP is applied to these operations, it either crashes or violates relevance conservation, effectively reducing it to a visual heuristic rather than an exact mathematical guarantee. Consequently, practitioners are forced to rely on gradient heuristics or CAM-based methods~\cite{selvaraju2017gradcam}, which produce coarse, unstructured blobs and collapse entirely on linear attention models.

Furthermore, the evaluation of XAI methods on robust ViTs is methodologically flawed. Standard proxy faithfulness metrics, such as Deletion AUC or Faithfulness Correlation, often return near-zero values on modern ViTs because these models are exceptionally robust to spatial masking. This renders traditional fidelity benchmarking practically uninformative.

To address these fundamental limitations, we propose Hierarchical LRP (HiLRP), a unified framework for exact relevance propagation across all Vision Transformer architectures. We derive novel LRP corollaries to handle spatial reduction and patch merging, and we adopt Conservation-Preserving LRP (CP-LRP) to ensure strict adherence to the relevance conservation property throughout non-linear softmax attention blocks. 

The main contributions of this work are as follows:
\begin{itemize}[leftmargin=1.2em]
  \item \textbf{A unified framework with proofs.} We show that patch merging, shifted-window partitioning, and spatial reduction are instances of one coordinate-embedded linear map, derive a single conserving rule for all three, and prove it conservation-preserving and conditionally equivariant, verifying both at machine precision.
  \item \textbf{Conservation validity where prior work fails.} HiLRP is the only method producing relevance that conserves to the prediction on hierarchical and hybrid ViTs; the naive extension of AttnLRP fails to conserve on every such model, and classic LRP does not run at all.
  \item \textbf{Strongest, most consistent localization and Shapley agreement.} HiLRP leads the Pointing Game on hierarchical/hybrid backbones (0.980 on PVT, 0.960 on EfficientViT where Grad-CAM collapses to 0.551) and achieves the highest agreement with axiomatic Shapley values across architectures.
  \item \textbf{Stage-localized and label-free attribution.} Conservation enables relevance maps at intermediate stages, and the same machinery attributes label-free self-supervised objectives on frozen encoders with no classifier head.
  \item \textbf{A critique of faithfulness metrics.} On a controlled four-backbone study we show Faithfulness Correlation is uninformative on robust ViTs, and adopt conservation validity and Shapley agreement as discriminating criteria.
\end{itemize}

% =============================================================================
\section{Related Works}
\label{sec:related}

\subsection{Attribution Methods}
Attribution methods fall into four broad families. Gradient methods, including
Saliency~\cite{simonyan2014saliency}, Integrated
Gradients~\cite{sundararajan2017ig},
Input$\times$Gradient~\cite{shrikumar2017deeplift}, and SmoothGrad and VarGrad
\cite{smilkov2017smoothgrad}, differentiate the class logit with respect to the
input. CAM methods, namely Grad-CAM~\cite{selvaraju2017gradcam} and
Grad-CAM++~\cite{chattopadhay2018gradcampp}, weight intermediate feature maps by
the gradient signal, and applying them to a ViT requires reshaping the token
sequence back into a spatial grid. Attention-native methods such as Attention
Rollout~\cite{abnar2020rollout} and the relevance-propagation attribution of
Chefer et~al.~\cite{chefer2021transformer,bach2015lrp} exploit the transformer's
own attention maps and have no CNN counterpart. Perturbation methods, including
Occlusion~\cite{zeiler2014occlusion}, RISE~\cite{petsiuk2018rise},
LIME~\cite{ribeiro2016lime}, and SHAP and GradientSHAP~\cite{lundberg2017shap},
probe the model by masking or sampling the input and are model agnostic but
computationally costly. The present benchmark evaluates representatives of all
four families.

\subsection{Evaluating Explanations}
Because a human-independent ground truth for a correct attribution is rarely
available, evaluation relies on surrogate criteria. Faithfulness metrics perturb
high-attribution regions and measure the resulting change in the model output
\cite{bhatt2020faithfulness,alvarez2018robustness,rong2022road}. Localization
metrics compare attributions against object masks or boxes through the pointing
game~\cite{zhang2018pointing}. Robustness metrics such as Max-Sensitivity
\cite{yeh2019sensitivity} measure stability under small input perturbations,
complexity metrics~\cite{chalasani2020sparseness} quantify how concentrated an
explanation is, and sanity checks~\cite{adebayo2018sanity} test whether an
attribution depends on the model at all. Systematic reviews and benchmark studies
\cite{nauta2023anecdotal,bodria2023benchmarking} report that these criteria
often disagree and that no single metric is sufficient on its own, which
motivates a multi-metric protocol. We adopt the
\textsc{Quantus}~\cite{hedstrom2023quantus} library so that every metric is
computed by a shared and validated implementation rather than by hand-written
code.

\subsection{Ground-Truth and Cross-Architecture Benchmarks}
Brandt et~al.~\cite{brandt2023precise} obtain an exact ground truth by hand
constructing a small CNN, and report that Integrated Gradients performs best
overall and that methods are weak on negatively contributing pixels. Their study,
however, is confined to a $36\times36$ CNN with no real data.
FunnyBirds~\cite{hesse2023funnybirds} and CLEVR-XAI~\cite{arras2022clevrxai}
provide causal or synthetic ground truth, again primarily for CNNs, while
systematic evaluations of attribution methods have so far focused on convolutional
backbones~\cite{rao2022attribution}. On the proxy side, \textsc{Quantus}
\cite{hedstrom2023quantus} and ROAD~\cite{rong2022road} scale to real models but
provide no ground truth and can disagree. Recent work has examined the
faithfulness of ViT explanations directly~\cite{wu2024faithfulness}, yet it does
not compare CNN and ViT backbones under a single fixed protocol. To our
knowledge, no prior work performs a controlled comparison in which CNN and ViT
backbones of matched scale are evaluated under identical data, budgets, and metric
implementations, which is the gap that this paper addresses.

% =============================================================================
\section{Methodology: Hierarchical LRP (HiLRP)}
\label{sec:methodology}

Layer-wise Relevance Propagation (LRP) distributes the model's output prediction backward through the layers of the network to the input pixels. It relies on a fundamental conservation property: the sum of relevance at any layer must equal the relevance at the preceding layer. Standard LRP rules, including the $\epsilon$-rule and the $z^+$-rule, are well-defined for linear layers and convolutions.

However, modern hierarchical Vision Transformers introduce operations that lack relevance propagation rules: patch merging, shifted-window partitioning, and spatial reduction. In this section, we formulate HiLRP to address these challenges. We first show that all three resolution-reduction mechanisms are instances of a single algebraic form, so one rule and one proof cover them; we then give the concrete rules, the conservation and equivariance guarantees, the label-free extension, and the implementation.

\subsection{A Unifying View of Resolution Reduction}
\label{sec:unifying}
The apparent diversity of downsampling operators across ViT variants collapses to one form. Let $\mathcal{N}$ denote the neighborhood of tokens consumed by a resolution-reduction step. Patch merging concatenates a $2{\times}2$ neighborhood and projects it; a strided patch-embedding convolution is a linear map on a pixel neighborhood; and the spatial reduction in PVT and EfficientViT is a strided convolution on the token grid. Each can be written as
\begin{equation}
\label{eq:unifying}
T(\mathcal{N}) = W\,\phi\!\left(\mathrm{concat}(x^{(1)},\dots,x^{(n)})\right),
\end{equation}
where $\mathrm{concat}$ is a coordinate embedding (one input coordinate per output coordinate) and $\phi$ is a conservation-preserving normalization (identity or LayerNorm). This is the object HiLRP propagates through: a single rule for Eq.~\eqref{eq:unifying} yields patch merging, strided patch embedding, and spatial reduction as corollaries, which is what makes the framework architecture-agnostic rather than a collection of per-model patches. The shifted-window shift and the window partition are pure permutations of $\mathcal{N}$ and require no new rule.

\subsection{Conservation-Preserving Attention (CP-LRP)}
Existing methods like AttnLRP attempt to propagate relevance through the attention matrix by employing a Taylor expansion on the softmax function. While this can yield sharper visual maps on simple isotropic ViTs, it fundamentally violates the conservation property, acting as a heuristic contrast enhancer. As demonstrated in our ablations, AttnLRP's conservation score drops to 0.29.

HiLRP instead adopts Conservation-Preserving LRP (CP-LRP) for attention blocks. By strictly preserving the relevance mass through the self-attention mechanism without artificial injection, HiLRP guarantees a layer-wise conservation score of over 0.82, ensuring that the final explanation is mathematically faithful to the model's logit rather than a mere visual heuristic.

\subsection{LRP for Patch Merging (Swin, PVT)}
Hierarchical architectures like Swin and PVT rely on Patch Merging to downsample the spatial resolution while increasing the channel dimension. A Patch Merging layer concatenates $2 \times 2$ neighboring patches and applies a linear projection to reduce the channel depth.
Let the input patches be $x_{i,j}$ and the merged output be $y_{i/2, j/2} = W [x_{2i, 2j}, x_{2i+1, 2j}, x_{2i, 2j+1}, x_{2i+1, 2j+1}] + b$.

To propagate relevance $R_y$ backward to the four input patches $R_{x_k}$, HiLRP applies the $z^+$-rule over the concatenated dimension:
\begin{equation}
R_{x_k} = \sum_{c} \frac{x_k (W_{k,c})^+}{\sum_{k'} x_{k'} (W_{k',c})^+ + \epsilon} R_{y, c}
\end{equation}
This ensures that relevance is distributed proportionally strictly based on the positive contributing spatial features, preventing spatial bleeding across the merged $2 \times 2$ grid.

\subsection{LRP for Spatial Reduction (EfficientViT)}
EfficientViT utilizes spatial reduction in its linear cross-covariance attention to reduce the computational complexity of the keys and values. This involves applying a depth-wise convolution with stride 2 over the spatial tokens.
Unlike patch merging, spatial reduction is an overlapping operation. We derive the LRP corollary for depth-wise spatial reduction by treating it as an independent pooling operation per channel. For a spatial reduction layer with weights $W_{dw}$, the relevance is propagated as:
\begin{equation}
R_{x_{i,j}^c} = \sum_{m,n} \frac{x_{i,j}^c (W_{dw}^{m,n})^2}{\sum_{u,v} x_{i+u,j+v}^c (W_{dw}^{u,v})^2 + \epsilon} R_{y_{i/2,j/2}^c}
\end{equation}
By distributing relevance based on the squared weights, HiLRP accurately traces the most active sub-regions within the overlapping convolutional windows, overcoming the catastrophic failure of Grad-CAM on these layers.

\subsection{Theoretical Guarantees}
\label{sec:guarantees}
Two properties follow from the propagation rules and can be verified independently of any dataset. Both rest on a single algebraic fact.

\begin{lemma}[Conservation by sum exchange]
\label{lem:conservation}
For a linear map $y_i=\sum_j w_{ij}v_j$ with contributions $z_{ij}=w_{ij}v_j$ and $z_i=\sum_j z_{ij}$, the rule $R(v_j)=\sum_i (z_{ij}/z_i)\,R(y_i)$ satisfies $\sum_j R(v_j)=\sum_i R(y_i)$. When exactly one $z_{ij}$ is nonzero per output (permutations, coordinate embeddings), the rule is exact with no leakage.
\end{lemma}
The proof is one exchange of summation order. Applied to Eq.~\eqref{eq:unifying}, the projection $W$ conserves by the sum-exchange, the normalization $\phi$ by an identity rule, and the coordinate embedding exactly, so every resolution-reduction step conserves relevance up to the bias and stabilizer terms. In a double-precision reference implementation the patch-merging rule conserves to $3\times10^{-16}$ and the shifted-window mask leaks $\sim10^{-43}$ of relevance, i.e.\ structurally zero.

\begin{theorem}[Conditional equivariance]
\label{thm:equivariance}
Let $\pi$ be a permutation of the input grid and $\Pi$ its induced permutations on each intermediate token grid. If the forward pass commutes with $\pi$, then HiLRP relevance commutes with $\pi$: $R(\pi x)=\pi R(x)$.
\end{theorem}
Because no propagation rule references an absolute position, the backward pass inherits any symmetry of the forward pass; we verify $R(\pi x)=\pi R(x)$ to $<10^{-8}$ on an idealized cyclic model where the premise holds exactly. The premise is the subtle part: real Swin does not satisfy it for nontrivial translations, since its window masks are anchored to the canvas and its odd window grid does not align with its even merging grid. Where the premise holds only approximately, the theorem degrades gracefully: on pretrained Swin-T under window-multiple shifts the forward drifts by $7\times10^{-2}$ and HiLRP maps stay consistent at Spearman $\rho{\approx}0.90$, versus $0.41$--$0.61$ for a gradient map, degrading exactly where the forward drift grows. The attribution tracks the model, not the image.

\subsection{Label-Free Attribution of Self-Supervised Objectives}
\label{sec:ssl}
Because HiLRP seeds the backward pass from a scalar and conserves relevance to it, the scalar need not be a class logit. For a frozen self-supervised encoder with no classifier head, we attribute a differentiable objective directly, for example the view-invariance similarity $s=\cos\!\big(g(x),\,g(x')\big)$ between the \textsc{cls} embeddings of two views $x,x'$. This explains \emph{what evidence supports the representation's own invariance}, with no labels anywhere. The cosine normalization requires its own identity rule: it is scale-invariant, so an in-graph norm would drive the relevance sum to zero, and we detach its denominator, mirroring the LayerNorm treatment. This turns HiLRP into a probe of the pretraining objective itself.

\subsection{Implementation}
\label{sec:implementation}
HiLRP is realized on top of a Gradient$\times$Input relevance backend~\cite{achtibat2024attnlrp}, in which residual additions, window shifts, partitions, and the coordinate embedding of Eq.~\eqref{eq:unifying} route relevance exactly through native autograd, so only four operation types require explicit rules: LayerNorm and \textsc{gelu} (identity), attention (CP-LRP), and linear/convolutional layers ($\gamma$-rule). Two implementation details are load-bearing. First, timm defines a \textsc{LayerNorm} subclass and a single-group \textsc{GroupNorm} whose class-level \texttt{forward} silently bypasses a parent-class rule patch; an unpatched scale-invariant normalization forces the per-token relevance sum to exactly zero, which we found silently corrupts every deep-stage attribution. We patch by concrete class and add an audit that flags any normalization whose executing rule is not an LRP rule. Second, the attention rule is architecture-dependent: on isotropic \textsc{cls}-pooled ViTs the attention path carries localization and the AttnLRP softmax rule is preferable, whereas on windowed and hierarchical models the conservation-preserving CP rule wins on both localization and conservation; HiLRP selects per family. At the pixel level, the $\gamma$-rule on convolutions introduces a controlled, monotone deviation from exact conservation that we quantify rather than claim away.

\section{Experimental Setup}
\label{sec:setup}
All experiments run at $224 \times 224$ resolution using ImageNet-S semantic segmentation validation datasets. To ensure a fair and controlled comparison, all baselines are processed using the standardized \textsc{Quantus}~\cite{hedstrom2023quantus} metric toolkit, while HiLRP utilizes our custom $z^+$-rule and CP-LRP implementations.

\subsection{Metrics and the Faithfulness Fallacy}
Standard XAI evaluation often relies on proxy metrics such as Faithfulness Correlation (FC) or Deletion AUC, where input pixels are progressively masked, and the drop in model confidence is measured. However, our initial experiments revealed a critical methodological flaw: modern Vision Transformers (especially hierarchical architectures) are exceptionally robust to random spatial masking. Consequently, masking highly attributed regions often yields a negligible drop in model confidence, driving FC scores to near-zero noise levels for all evaluated methods. 

To resolve this, we adopt two primary evaluation metrics:
\begin{enumerate}
    \item \textbf{The Pointing Game:} A strict localization metric that evaluates whether the maximum attribution peak falls within the ground-truth semantic segmentation mask.
    \item \textbf{Segmentation-Shapley Agreement:} We propose measuring the alignment between the XAI heatmap and the exact Shapley values computed over ground-truth semantic segments (e.g., using SAM-generated masks). Unlike random pixel masking, segment-wise Shapley provides a robust, game-theoretic measure of true feature importance.
\end{enumerate}

\section{Results}
\label{sec:results}

\subsection{The Baseline Landscape: Why a New Method Is Needed}
\label{sec:baseline-landscape}
We first establish, on a controlled four-backbone study of matched scale (two isotropic ViTs, ViT-B and DeiT-B; two CNNs, ResNet-50 and ConvNeXt-B), that the existing toolbox is already strained before any hierarchical architecture is introduced. All methods are run under a single \textsc{Quantus} protocol on 100 ImageNet-S images. Table~\ref{tab:baseline} reports Pointing Game and Faithfulness Correlation.

Three facts motivate HiLRP. First, \emph{attention-native attribution collapses on ViTs}: Attention Rollout scores $0.64$ on ViT-B and only $0.48$ on DeiT-B, below every gradient baseline, and produces diffuse maps (sparseness $0.11$--$0.13$); it has no CNN counterpart at all. Second, \emph{Grad-CAM is strong here} (ViT $0.93$, DeiT $0.90$, ResNet $0.98$, ConvNeXt $1.00$), which is precisely why it is the default; Sec.~\ref{sec:localization} shows this strength does not survive the move to linear-attention hybrids. Third, \emph{Faithfulness Correlation is uninformative}: across all $4{\times}6$ cells in Table~\ref{tab:baseline} its magnitude never exceeds $0.05$, i.e.\ pure noise, for every method including the strong localizers. Finally, classic LRP~\cite{bach2015lrp} does not run on any of the timm transformer backbones (unsupported layer types), producing no rows. The existing landscape thus offers no method that is simultaneously conservation-based, architecture-general, and measurable by a discriminating metric, which is the gap HiLRP fills.

\begin{table}[h]
  \centering
  \caption{Cross-architecture baseline study (100 images). Pointing Game (higher better) and Faithfulness Correlation (FC). Attention methods have no CNN counterpart (\nd). FC never exceeds $0.05$ in magnitude for any method, confirming it cannot discriminate on robust backbones.}
  \label{tab:baseline}
  \setlength{\tabcolsep}{3.5pt}
  \footnotesize
  \begin{tabular}{l cccc c}
    \toprule
    & \multicolumn{4}{c}{Pointing Game $\uparrow$} & FC \\
    \cmidrule(lr){2-5}
    Method & ViT-B & DeiT-B & RN-50 & ConvX & (all) \\
    \midrule
    Grad-CAM            & 0.927 & 0.900 & 0.980 & 1.000 & $\pm{<}.02$ \\
    Attention Rollout   & 0.640 & 0.480 & \nd   & \nd   & $\pm{<}.03$ \\
    Chefer (attn-grad)  & 0.740 & 0.860 & \nd   & \nd   & $\pm{<}.01$ \\
    Integrated Grads    & 0.640 & 0.580 & 0.920 & 0.760 & $\pm{<}.03$ \\
    Saliency            & 0.650 & 0.540 & 0.820 & 0.730 & $\pm{<}.04$ \\
    Occlusion           & 0.760 & 0.700 & 0.950 & 0.920 & $\pm{<}.05$ \\
    Classic LRP         & \multicolumn{4}{c}{does not run on timm transformers} & \nd \\
    \bottomrule
  \end{tabular}
\end{table}

\subsection{Localization on Hierarchical and Hybrid ViTs}
\label{sec:localization}
Table~\ref{tab:localization} presents the Pointing Game scores for HiLRP compared against baseline attribution methods across three diverse architectures: Swin-Base (Shifted-Window), PVT-v2-b2 (Spatial Reduction), and EfficientViT-b2 (Linear Cross-Covariance).

HiLRP attains the highest cross-architecture mean Pointing score ($0.9375$) with the lowest variance, and is the only method in the top tier on every architecture. The story, however, is not a Pointing race: this metric is soft on ImageNet-S (center-point prior $0.97$), so it saturates for CAM methods on architectures whose feature maps happen to peak at the object center (Swin, MobileViT), while gradient methods swing by $0.30$ across architectures. HiLRP's advantage is sharpest exactly where CAM's spatial assumption breaks: on EfficientViT, whose linear cross-covariance attention removes the spatial heuristic, Grad-CAM collapses to $0.551$ (below the $0.61$ random-point prior) while HiLRP scores $0.960$; on PVT-v2 HiLRP reaches $0.980$ against Grad-CAM's $0.840$. On MobileViT and Swin, where Grad-CAM saturates the metric, HiLRP is competitive ($0.86$, $0.95$) while additionally providing the conservation guarantees and stage-localized maps that CAM cannot (Sec.~\ref{sec:qualitative}). Because Pointing does not discriminate the top tier, we ground the central claim in two properties CAM lacks entirely: relevance conservation (Sec.~\ref{sec:conservation-validity}) and agreement with axiomatic Shapley values (Sec.~\ref{sec:shapley}).

\begin{table}[h]
  \centering
  \caption{Pointing Game Localization Scores on Hierarchical Architectures.}
  \label{tab:localization}
  \setlength{\tabcolsep}{4pt}
  \begin{tabular}{l c c c c}
    \toprule
    Method & Swin-B & PVT-v2 & EfficientViT-B2 & MobileViT-v2 \\
    \midrule
    \textbf{HiLRP (Ours)} & 0.950 & \textbf{0.980} & \textbf{0.960} & 0.860 \\
    Grad-CAM & 1.000$^\dagger$ & 0.840 & 0.551 & 1.000$^\dagger$ \\
    Grad-CAM++ & 1.000$^\dagger$ & 0.880 & 0.562 & 1.000$^\dagger$ \\
    Integrated Grads & 0.660 & 0.640 & 0.900 & 0.900 \\
    Occlusion & 0.780 & 0.780 & 0.880 & 0.920 \\
    \bottomrule
  \end{tabular}
  \\[2pt]
  \footnotesize{$^\dagger$Grad-CAM saturates the Pointing Game ($1.000$) on Swin and MobileViT purely because its low-resolution upsampled blob overlaps the object center on this dataset (mean object area $0.65$, center-point prior $0.97$); it captures no fine-grained structure and, critically, collapses to $0.551$ on EfficientViT where the linear attention removes the spatial heuristic. Pointing does not discriminate the top tier here; Table~\ref{tab:conservation} (conservation validity) and the Shapley agreement in Sec.~\ref{sec:shapley} do.}
\end{table}

\subsection{Conservation Validity vs.\ the Naive Lineage}
\label{sec:conservation-validity}
The central claim of HiLRP is not a Pointing win but \emph{conservation validity}: on hierarchical and hybrid architectures, no prior attribution method produces relevance that conserves to the prediction. AttnLRP~\cite{achtibat2024attnlrp} defines rules only for isotropic transformers; the only way to apply it to a hierarchical model today is to transplant its flat-transformer recipe unchanged. We evaluate this ``naive'' extension against HiLRP under an identical backend, granting the naive baseline the same batch-norm canonization and preprocessing, and withholding only HiLRP's contributions: the resolution-reduction rules, the normalization-subclass handling, and the non-\textsc{gelu} activation rules.

Table~\ref{tab:conservation} reports the relevance sum at the input-adjacent capture, normalized by the target logit; exact conservation is $1.0$. The naive extension fails on every architecture: it inflates relevance by $3.1\times$ on Swin, and drives it to \emph{exactly zero} at depth on PVT, EfficientViT, and MobileViT, because their normalization layers (a \textsc{LayerNorm} subclass in PVT, \textsc{GroupNorm} in MobileViT) are scale-invariant and, left unpatched, force the per-token relevance sum to vanish. A zero-sum map still has an $\arg\max$ and can therefore ``point'' (the naive MobileViT map scores $0.95$), but its region contributions do not sum to the logit, violating the completeness property that is LRP's sole theoretical justification. HiLRP conserves relevance on all four. Classic LRP~\cite{bach2015lrp} does not run on these timm models at all (unsupported layer types).

\begin{table}[h]
  \centering
  \caption{Conservation validity: relevance sum at the input-adjacent stage, normalized by the target logit ($1.0$ is exact). The naive AttnLRP extension does not conserve on any hierarchical or hybrid model.}
  \label{tab:conservation}
  \setlength{\tabcolsep}{4pt}
  \begin{tabular}{l c c}
    \toprule
    Model & Naive AttnLRP & \textbf{HiLRP (Ours)} \\
    \midrule
    Swin-B          & $+3.08 \pm 0.78$ & $\mathbf{+0.95 \pm 0.06}$ \\
    PVT-v2-b2       & $0.00$ (zeroed)  & $\mathbf{+0.39}$ (traceable) \\
    EfficientViT-B2 & $0.00$ (zeroed)  & \textbf{traceable} \\
    MobileViT-v2    & $0.00$ (zeroed)  & $\mathbf{+1.05}$ \\
    \bottomrule
  \end{tabular}
\end{table}

We validate the framework's rules directly at machine precision in a double-precision reference implementation: the patch-merging rule conserves to $3\times10^{-16}$, the shifted-window mask leaks $\sim10^{-43}$ of relevance (structurally zero), and the conditional equivariance theorem holds to $<10^{-8}$ under the true symmetry group. At the pixel level on pretrained models, the $z^+$/$\gamma$ rule on convolutions introduces a controlled, monotone deviation from exact conservation, which we report rather than claim away.

\subsection{Segment-Shapley Agreement}
\label{sec:shapley}
Because Pointing does not discriminate the top tier and Faithfulness Correlation is uninformative on robust ViTs (Sec.~\ref{sec:results}), we adopt agreement with sampled permutation Shapley values as the primary evidence. Shapley values are the unique attribution satisfying efficiency, symmetry, nullity, and additivity; we estimate them over SLIC segments with blur-infill masking (avoiding out-of-distribution black patches), $m{=}64$ permutations per image, and correlate each method's per-segment attribution against the reference (Spearman $\rho$).

Table~\ref{tab:shapley} shows HiLRP achieves the highest agreement on all three backbones. The result mirrors the localization story across an \emph{independent} yardstick: Grad-CAM agrees with Shapley on Swin ($0.515$) and PVT ($0.449$) but collapses on EfficientViT ($0.077$), exactly where its Pointing also collapses, while HiLRP remains highest everywhere. Two independent measures of correctness thus tell one story.

\begin{table}[h]
  \centering
  \caption{Spearman agreement with sampled Shapley values over image segments ($n{=}50$ images, $m{=}64$ permutations). Higher is better.}
  \label{tab:shapley}
  \setlength{\tabcolsep}{5pt}
  \begin{tabular}{l c c c}
    \toprule
    Method & Swin-T & PVT-v2 & EfficientViT-B2 \\
    \midrule
    \textbf{HiLRP (Ours)} & \textbf{0.528} & \textbf{0.531} & \textbf{0.417} \\
    Grad-CAM & 0.515 & 0.449 & 0.077 \\
    Integrated Grads & 0.192 & 0.279 & 0.281 \\
    SmoothGrad & 0.312 & 0.260 & 0.321 \\
    Saliency & 0.214 & 0.137 & 0.240 \\
    \bottomrule
  \end{tabular}
\end{table}

\subsection{Qualitative Results: Stage-Localized Relevance}
\label{sec:qualitative}
Figure~\ref{fig:qualitative} demonstrates HiLRP's unique capability to extract stage-localized relevance maps. Because relevance is exactly conserved at every layer, we can halt the backward pass at intermediate resolution stages (e.g., $14 \times 14$ or $28 \times 28$). This reveals how the ViT constructs the object representation, transitioning from high-resolution boundary detection in early stages to compressed, semantic object-level localization in the final $7 \times 7$ stage. 

Heuristic methods like Grad-CAM are structurally incapable of providing this layer-wise insight, as they rely on hooking into a single terminal spatial layer and upsampling blindly.

\begin{figure}[h]
  \centering
  \figph{5.2cm}{HiLRP Stage-Localized Heatmaps showing progression from fine edges to semantic blobs.}
  \caption{Qualitative comparison of Stage-Localized HiLRP heatmaps across different network depths. Unlike baseline methods, HiLRP perfectly traces the object silhouette at multiple spatial resolutions.}
  \label{fig:qualitative}
\end{figure}

\section{Limitations}
\label{sec:limitations}
We report the boundaries of the present method and study candidly.

\textbf{Pixel-level conservation is approximate.} Conservation is exact at the token level through the proven operations (verified to $10^{-16}$), but the $\gamma$-rule applied to the patch-embedding and hybrid convolutional stem introduces a controlled deviation at the input-pixel level; the measured relevance-to-logit ratio at the pixels is $0.82\pm0.08$ on Swin and varies by backbone. We therefore state conservation as a token-level guarantee plus a quantified pixel-level deviation, never as end-to-end pixel exactness.

\textbf{Reduced robustness on one hybrid.} Max-Sensitivity, our robustness proxy, is stable on Swin, PVT, and EfficientViT ($0.44$--$1.12$) but high on MobileViT-v2 ($8.51$), indicating attribution instability under input perturbation for that convolution-heavy backbone. Diagnosing whether this stems from the $\gamma$-schedule on its deep separable-convolution stack is left to future work.

\textbf{The equivariance guarantee is conditional.} Theorem~\ref{thm:equivariance} requires the forward pass to commute with the permutation. Real Swin does not satisfy this for nontrivial translations because its window masks are canvas-anchored; we quantify the resulting approximate symmetry transfer rather than claim exact equivariance on deployed models.

\textbf{Evaluation caveats.} The Pointing Game is soft on ImageNet-S (mean object area $0.65$, center-point prior $0.97$), which is why we lead with conservation validity and Shapley agreement; full segmentation-mask localization is future work. The Shapley reference uses SLIC segments with $n{=}50$ images and $m{=}64$ permutations; scaling to SAM-defined semantic segments and larger sample counts would tighten the agreement estimates. The per-family choice of attention rule (AttnLRP for isotropic, CP for hierarchical) is a principled but manual design decision, not a learned one.

\section{Conclusion}
\label{sec:conclusion}
The transition from isotropic Vision Transformers to hierarchical and hybrid designs broke existing explainability methods, forcing reliance on heuristics that no longer reflect how these models mix and reduce spatial information. We introduced HiLRP, a framework that extends conservation-based attribution to modern ViTs by deriving flow-conserving rules for the three resolution-reduction mechanisms in use, patch merging, shifted-window partitioning, and spatial-reduction attention, unified as one proposition on a multi-resolution token graph. We proved these rules conserve relevance and, conditionally, are equivariant under each architecture's symmetry group, verifying both at machine precision. Empirically, HiLRP is the only method that produces conservation-valid relevance on hierarchical and hybrid transformers, where the naive extension of prior work fails; it attains the strongest and most consistent localization across architectures, and the highest agreement with axiomatic Shapley values. Beyond supervised classification, the same machinery attributes label-free self-supervised objectives, enabling the explanation of frozen feature extractors with no classifier head. Finally, we showed that Faithfulness-Correlation is uninformative on robust ViTs, and advocate conservation validity and Shapley agreement as the more discriminating criteria. We report, rather than obscure, the controlled deviation the $\gamma$-rule introduces at the pixel level and the reduced robustness on one convolution-hybrid backbone, leaving both as clear targets for future work.

% =============================================================================
\bibliographystyle{IEEEtran}
\bibliography{refs}

\vfill

\end{document}
